% GENERATED FILE. Edit canonical lessons, then run npm run generate:books.
% canonical-source-hash: fnv1a64:4f5d89d0d86c9c7e
% canonical-lessons: PT-C455-levantamento, PT-C455-cobrar, PT-C455-debitar, PT-C455-extrato, PT-C455-comissao

\chapter{No banco}
\label{ch:pt-no-banco}
\hlchaptermodality{455}

\begin{chapteropening}
\textbf{By the end of this chapter:} \emph{I can manage money at the bank: a withdrawal, charging, debiting, the statement and the fee.}

Complete the last lesson of chapter 455: I can manage money at the bank: a withdrawal, charging, debiting, the statement and the fee.
\end{chapteropening}

\section[o levantamento]{o levantamento — the cash withdrawal; the survey}

\label{lesson:PT-C455-levantamento}



\begin{quote}
Before the new one: say the Portuguese for to dispatch, then the Portuguese for to register.
\end{quote}

\subsection*{You'll want to know: o levantamento}
\textbf{o levantamento} — \textquotedblleft{}the cash withdrawal; the survey\textquotedblright{}.

\begin{quote}
Fiz um levantamento no multibanco.
\end{quote}

\begin{quote}
\emph{I made a withdrawal at the cash machine.}
\end{quote}

\emph{Levantar dinheiro} — to withdraw cash.

\begin{cousinweb}
\emph{levantamento} is built on \emph{levantar} 'to lift, to raise', which goes back to Latin \emph{levāre} 'to lighten, to raise' (Lewis \& Short).
\end{cousinweb}

\begin{grammarlens}[title={What you've built}]
One more word for telling, explaining and describing at B1.
\end{grammarlens}

\subsection*{Your turn}
\begin{itemize}
\raggedright
  \item \emph{Say it:} \emph{o levantamento}
  \item \emph{Say it:} \emph{o levantamento}, once more
  \item \emph{Say it:} a sentence of your own with \emph{o levantamento}
  \item \emph{Recall:} say the Portuguese for to dispatch, then the Portuguese for to register, then say \emph{o levantamento} again
\end{itemize}

\begin{tcolorbox}[breakable,colback=teal!4,colframe=teal!35!black,title={Before you move on}]
What does o levantamento mean? (\textquotedblleft{}The cash withdrawal; the survey\textquotedblright{}.) Say it once more.
\end{tcolorbox}
\section[cobrar]{cobrar — to charge; to collect (a debt)}

\label{lesson:PT-C455-cobrar}



\begin{quote}
Before the new one: say the Portuguese for to register, then the Portuguese for the cash withdrawal.
\end{quote}

\subsection*{You'll want to know: cobrar}
\textbf{cobrar} — \textquotedblleft{}to charge; to collect (a debt)\textquotedblright{}.

\begin{quote}
O banco cobra dois euros por cada transferência.
\end{quote}

\begin{quote}
\emph{The bank charges two euros for each transfer.}
\end{quote}

\emph{Cobrar caro} — to charge a lot.

\begin{cousinweb}
\emph{cobrar} is a shortening of \emph{recobrar} 'to recover' (EtymDB), which is inherited from Latin \emph{recuperāre} 'to get back'.
\end{cousinweb}

\begin{grammarlens}[title={What you've built}]
One more word for telling, explaining and describing at B1.
\end{grammarlens}

\subsection*{Your turn}
\begin{itemize}
\raggedright
  \item \emph{Say it:} \emph{cobrar}
  \item \emph{Say it:} \emph{cobrar}, once more
  \item \emph{Say it:} a sentence of your own with \emph{cobrar}
  \item \emph{Recall:} say the Portuguese for to register, then the Portuguese for the cash withdrawal, then say \emph{cobrar} again
\end{itemize}

\begin{tcolorbox}[breakable,colback=teal!4,colframe=teal!35!black,title={Before you move on}]
What does cobrar mean? (\textquotedblleft{}To charge; to collect (a debt)\textquotedblright{}.) Say it once more.
\end{tcolorbox}
\section[debitar]{debitar — to debit, to take (money) from an account}

\label{lesson:PT-C455-debitar}



\begin{quote}
Before the new one: say the Portuguese for the cash withdrawal, then the Portuguese for to charge.
\end{quote}

\subsection*{You'll want to know: debitar}
\textbf{debitar} — \textquotedblleft{}to debit, to take (money) from an account\textquotedblright{}.

\begin{quote}
O banco debitou a renda no dia um.
\end{quote}

\begin{quote}
\emph{The bank took the rent out on the first.}
\end{quote}

\emph{O débito direto} — a direct debit.

\begin{cousinweb}
\emph{debitar} is built on \emph{débito} 'debit, debt', a learned borrowing from Latin \emph{dēbitum} 'what is owed', from \emph{dēbēre} 'to owe' (Lewis \& Short).
\end{cousinweb}

\begin{grammarlens}[title={What you've built}]
One more word for telling, explaining and describing at B1.
\end{grammarlens}

\subsection*{Your turn}
\begin{itemize}
\raggedright
  \item \emph{Say it:} \emph{debitar}
  \item \emph{Say it:} \emph{debitar}, once more
  \item \emph{Say it:} a sentence of your own with \emph{debitar}
  \item \emph{Recall:} say the Portuguese for the cash withdrawal, then the Portuguese for to charge, then say \emph{debitar} again
\end{itemize}

\begin{tcolorbox}[breakable,colback=teal!4,colframe=teal!35!black,title={Before you move on}]
What does debitar mean? (\textquotedblleft{}To debit, to take (money) from an account\textquotedblright{}.) Say it once more.
\end{tcolorbox}
\section[o extrato]{o extrato — the bank statement; the extract}

\label{lesson:PT-C455-extrato}



\begin{quote}
Before the new one: say the Portuguese for to charge, then the Portuguese for to debit.
\end{quote}

\subsection*{You'll want to know: o extrato}
\textbf{o extrato} — \textquotedblleft{}the bank statement; the extract\textquotedblright{}.

\begin{quote}
Pedi um extrato da conta no balcão.
\end{quote}

\begin{quote}
\emph{I asked for an account statement at the counter.}
\end{quote}

\emph{Extrato de tomate} — tomato purée.

\begin{cousinweb}
\emph{extrato} is a learned borrowing from Latin \emph{extractum} 'something drawn out' (EtymDB), from \emph{extrahere} 'to draw out', on \emph{trahere} 'to pull'.
\end{cousinweb}

\begin{grammarlens}[title={What you've built}]
One more word for telling, explaining and describing at B1.
\end{grammarlens}

\subsection*{Your turn}
\begin{itemize}
\raggedright
  \item \emph{Say it:} \emph{o extrato}
  \item \emph{Say it:} \emph{o extrato}, once more
  \item \emph{Say it:} a sentence of your own with \emph{o extrato}
  \item \emph{Recall:} say the Portuguese for to charge, then the Portuguese for to debit, then say \emph{o extrato} again
\end{itemize}

\begin{tcolorbox}[breakable,colback=teal!4,colframe=teal!35!black,title={Before you move on}]
What does o extrato mean? (\textquotedblleft{}The bank statement; the extract\textquotedblright{}.) Say it once more.
\end{tcolorbox}
\section[a comissão]{a comissão — the fee, the commission; the committee}

\label{lesson:PT-C455-comissao}



\begin{quote}
Before the new one: say the Portuguese for to debit, then the Portuguese for the bank statement.
\end{quote}

\subsection*{You'll want to know: a comissão}
\textbf{a comissão} — \textquotedblleft{}the fee, the commission; the committee\textquotedblright{}.

\begin{quote}
Este cartão não tem comissão anual.
\end{quote}

\begin{quote}
\emph{This card has no annual fee.}
\end{quote}

\emph{A comissão de moradores} — the residents' committee.

\begin{cousinweb}
\emph{comissão} is a learned borrowing from Latin \emph{commissiō} 'bringing together; commission' (EtymDB), from \emph{committere}, on \emph{mittere} 'to send'.
\end{cousinweb}

\begin{grammarlens}[title={What you've built}]
One more word for telling, explaining and describing at B1.
\end{grammarlens}

\subsection*{Your turn}
\begin{itemize}
\raggedright
  \item \emph{Say it:} \emph{a comissão}
  \item \emph{Say it:} \emph{a comissão}, once more
  \item \emph{Say it:} a sentence of your own with \emph{a comissão}
  \item \emph{Recall:} say the Portuguese for to debit, then the Portuguese for the bank statement, then say \emph{a comissão} again
\end{itemize}

\begin{tcolorbox}[breakable,colback=teal!4,colframe=teal!35!black,title={Before you move on}]
What does a comissão mean? (\textquotedblleft{}The fee, the commission; the committee\textquotedblright{}.) Say it once more.
\end{tcolorbox}
